\documentclass[9pt,landscape]{extarticle}
\usepackage[T1]{fontenc}
\usepackage{lmodern}
\usepackage[a4paper,landscape,margin=1cm,top=0.9cm,bottom=0.9cm]{geometry}
\usepackage{multicol}
\usepackage[table]{xcolor}
\usepackage{titlesec}
\usepackage{array}
\usepackage{tabularx}
\usepackage{enumitem}
\usepackage{fancyhdr}

% onedark "warmer"-ish palette
\definecolor{accent}{HTML}{E06C75}
\definecolor{accent2}{HTML}{61AFEF}
\definecolor{keybg}{HTML}{EDEFF2}
\definecolor{keyfg}{HTML}{2C323C}
\definecolor{muted}{HTML}{5C6370}

\setlength{\parindent}{0pt}
\setlength{\columnsep}{0.7cm}
\setlength{\tabcolsep}{3pt}

\titleformat{\section}{\normalfont\bfseries\color{accent}}{}{0em}{}[\vspace{-0.9em}\color{accent}\rule{\linewidth}{0.6pt}]
\titlespacing{\section}{0pt}{1.1ex plus .3ex}{0.7ex}

\newcommand{\key}[1]{\colorbox{keybg}{\textcolor{keyfg}{\ttfamily\footnotesize #1}}}
% \tabularx/\endtabularx, not begin/end: tabularx scans for a literal \end{tabularx}.
\newenvironment{keys}
  {\par\nobreak\tabularx{\linewidth}{@{}l@{\hspace{5pt}}X@{}}}
  {\endtabularx\par\vspace{0.3ex}}
\newcommand{\note}[1]{\par\textcolor{muted}{\itshape\footnotesize #1}\par\vspace{0.4ex}}
\newcommand{\subhead}[1]{\par\vspace{0.3ex}\textbf{\footnotesize #1}\nobreak}

\pagestyle{fancy}
\fancyhf{}
\renewcommand{\headrulewidth}{0pt}
\fancyfoot[R]{\textcolor{muted}{\footnotesize\thepage}}

\begin{document}

{\LARGE\bfseries Neovim Cheat Sheet}\quad
{\color{muted}\small nixvim-config \textbf{---} leader is \key{<Space>}\quad modes: \textsc{n}ormal \textsc{i}nsert \textsc{v}isual}
\vspace{0.4ex}
\hrule height 1pt
\vspace{1ex}

\begin{multicols}{3}

\section{Editing}
\begin{keys}
\key{jj}            & \textsc{i} leave insert mode \\
\key{<Space>}       & \textsc{v} leave visual mode \\
\key{<}             & \textsc{v} outdent, keep selection \\
\key{>}             & \textsc{v} indent, keep selection \\
\end{keys}

\section{Buffers}
\begin{keys}
\key{<c-w>n}        & next buffer \\
\key{<c-w>p}        & previous buffer \\
\key{<c-w>y}        & close buffer, keep window layout \\
\key{<leader>b}     & list / switch buffers (Telescope) \\
\end{keys}
\note{Tabs on top come from bufferline; \key{<c-w>y} uses bufdelete so the split survives. \texttt{:Bwipeout} does the same from the cmdline. It is a no-op in unlisted buffers (neo-tree, help).}

\section{File tree (neo-tree)}
\begin{keys}
\key{<leader>t}     & toggle tree \\
\key{<leader>T}     & reveal current file in tree \\
\key{l}             & expand dir / open file \\
\key{h}             & collapse dir / go to parent \\
\key{<cr>}          & toggle dir / open file \\
\key{P}             & preview file \\
\key{e}             & toggle fit-to-content width \\
\key{a d r c m}     & add, delete, rename, copy, move \\
\key{H}             & toggle hidden files \\
\key{?}             & all tree mappings \\
\end{keys}
\note{Width 40. \key{<Space>} is unbound in the tree so leader combos fire without waiting out timeoutlen. Dotfiles are shown; the tree follows the current file.}

\section{Find \& search (Telescope)}
\begin{keys}
\key{<leader>f}     & find project files \\
\key{<leader>g}     & live grep (root dir) \\
\key{<leader>b}     & buffers \\
\key{<leader>cb}    & fuzzy find in current buffer \\
\key{gr}            & find usages (LSP references) \\
\end{keys}
\subhead{Inside a picker}
\begin{keys}
\key{<c-j>} / \key{<c-k>} & next / previous result \\
\key{<cr>}          & open \\
\key{<c-x>}         & open in split \\
\key{<c-v>}         & open in vsplit \\
\key{<c-t>}         & open in new tab \\
\key{<esc>} / \key{<c-e>} & close (any mode) \\
\end{keys}
\note{\key{<esc>} closes rather than dropping to normal mode --- use \key{jj} if you want normal mode in the prompt.}

\columnbreak

\section{LSP}
\begin{keys}
\key{K}             & hover docs \\
\key{gd}            & go to definition \\
\key{gD}            & go to declaration \\
\key{gi}            & go to implementation \\
\key{gt}            & go to type definition \\
\key{gr}            & references / call sites \\
\key{<leader>rn}    & rename symbol \\
\key{<leader>ca}    & code action \\
\key{<leader>sh}    & signature help \\
\key{<leader>wa}    & add workspace folder \\
\key{<leader>wr}    & remove workspace folder \\
\end{keys}

\section{Diagnostics}
\begin{keys}
\key{<leader>j}     & next diagnostic \\
\key{<leader>k}     & previous diagnostic \\
\key{<leader>e}     & show diagnostic in float \\
\end{keys}

\section{Completion (blink.cmp)}
\begin{keys}
\key{<c-n>}         & open menu / toggle docs \\
\key{<c-j>}         & select next item \\
\key{<c-k>}         & select previous item \\
\key{<tab>}         & accept completion \\
\key{<c-e>}         & hide menu \\
\end{keys}
\note{Sources: LSP, path, snippets, buffer, and Copilot (full profile) ranked to the top. Signature help shows while typing arguments.}

\section{Copilot Chat}
\begin{keys}
\key{<leader>cc}    & open chat (\textsc{n} and \textsc{v}) \\
\key{<c-space>}     & submit prompt \\
\key{<c-r>}         & reset conversation \\
\key{<c-e>}         & close chat \\
\key{gd}            & show diff of a suggestion \\
\end{keys}
\note{Opens in a vertical split on the right, in insert mode. Buffers, selection and glob are available as resources (\texttt{\#buffers}, \texttt{\#selection}). Full profile only; run \texttt{:Copilot auth} once.}

\section{Git}
\begin{keys}
\key{]c} / \key{[c} & next / previous hunk (gitsigns) \\
\end{keys}
\begin{keys}
\texttt{:DiffviewOpen}        & side-by-side diff tab \\
\texttt{:DiffviewFileHistory} & history of a file \\
\texttt{:DiffviewClose}       & close the diff tab \\
\end{keys}
\note{In diffview: \key{<tab>} next file, \key{<s-tab>} previous, \key{g?} help.}

\columnbreak

\section{Code review (review.nvim)}
\begin{keys}
\key{<leader>rr}    & review working tree \\
\key{<leader>rc}    & review commits (picker) \\
\key{<leader>rb}    & review branch against main \\
\key{<leader>ra}    & note on line / selection, any file \\
\key{<leader>re} / \key{<leader>rd} & edit / delete note at cursor \\
\key{<leader>rx}    & export comments to clipboard \\
\end{keys}
\subhead{Inside the review}
\begin{keys}
\key{i}             & add comment (\textsc{v}: range) \\
\key{e} / \key{d}   & edit / delete comment \\
\key{]n} / \key{[n} & next / previous comment \\
\key{c} / \key{C}   & list / export comments \\
\key{q}             & close: export, then archive \\
\end{keys}

\section{Windows \& splits (vim defaults)}
\begin{keys}
\key{<c-w>s} / \key{<c-w>v} & split / vsplit \\
\key{<c-w>hjkl}     & move between windows \\
\key{<c-w>HJKL}     & move the window itself \\
\key{<c-w>=}        & equalize sizes \\
\key{<c-w>o}        & only this window \\
\key{<c-w>q}        & quit window \\
\end{keys}
\note{\texttt{splitright} is on, so vsplits open to the right.}

\section{Other commands}
\begin{keys}
\texttt{:MarkdownPreview} & preview markdown in browser \\
\texttt{:TodoTelescope}   & list \texttt{TODO}/\texttt{FIX}/\texttt{HACK} comments \\
\texttt{:checkhealth}     & diagnose the setup \\
\end{keys}
\note{Highlighted comment tags: FIX, TODO, HACK, WARN, PERF, XXX, INFO, NOTE, TEST.}

\section{Worth remembering (vim defaults)}
\begin{keys}
\key{gv}            & reselect last visual block \\
\key{<c-o>} / \key{<c-i>} & back / forward in jumplist \\
\key{<c-\^{}}         & alternate (previous) file \\
\key{ciw} / \key{diw}     & change / delete word under cursor \\
\key{.}             & repeat last change \\
\key{\%}             & jump to matching bracket \\
\key{zz} / \key{zt}       & center / top the cursor line \\
\key{gg=G}          & reindent the whole file \\
\key{q\{a-z\}} \dots\ \key{q} & record macro, \key{@a} replays \\
\end{keys}

\section{Profiles}
\note{The light profile ships only \texttt{nil\_ls} (Nix). The full profile adds ltex, pyright, clangd, fortls, jsonls, yamlls, jdtls, rust-analyzer, plus Copilot and markdown-preview --- LSP keys do nothing without a server for the filetype.}

\end{multicols}
\end{document}
